\documentstyle[%


righttag%


,11pt%


,amssymb%


,russian%               remove in Eng.


,izvru%                 Set izven% in Eng.


]{amsart}





% 





\newcommand{\ve}{\varepsilon}


\newcommand{\fy}{\infty}


\newcommand{\vt}{\vartheta}


\newcommand{\R}{\Bbb R}


\newcommand{\A}{\Bbb A}


\newcommand{\col}{\operatorname{col}}


\newcommand{\rank}{\operatorname{rank}}


\newcommand{\Int}{\operatorname{int}}


\newcommand{\Vc}{\operatorname{vec}}


\newcommand{\Spp}{\operatorname{Sp}}


\newcommand{\tts}[1]{}





%\hoffset=-15mm%                  For Eng.


 


\setcounter{page}{45}


\god{1999}


\nomer{2~(441)} %                        Remove in Eng.


%\nomer{Vol.\,43, No.\,2}%               Set in Eng.


%\str{\pageref{begin}--\pageref{end}}%  Set in Eng.


\udk{517.934}


\author{\it В.А.\,Зайцев, Е.Л.\,Тонков}


% NB:   ^^ remove in Eng.


\title{Достижимость, согласованность и метод поворотов \\


В.М.\,Миллионщикова}





\thanks{Работа поддержана Российским фондом


фундаментальных исследований (грант 97--01--00413) и конкурсным центром


Удмуртского госуниверситета (грант 97--04).}


\date{}


%\pagestyle{myheadings}%         Set in Eng.





\pagestyle{plain} %              Remove in Eng.


\begin{document}





%\thispagestyle{plain}%          Set in Eng.





\maketitle


\label{begin}





Здесь исследуются условия, при которых существует принципиальная


возможность построения управления, создающего


``хорошее окружение'' тривиальному решению уравнения


\begin{equation}


\label{01}


\Dot x=v_0(t,x)+u_1v_1(t,x)+\dotsb+u_ rv_r(t,x),\quad (t,x)\in\R^{1+n},


\end{equation}


где $u\doteq(u_1,\dotsc,u_r)\in U$, $0\in U$. Пусть, например, требуется


построить такое управление $u(t)\in U$, что уравнение (\ref{01})


структурно устойчиво (в окрестности нуля). Оказывается, что эта задача


разрешима, если линейное уравнение


\begin{equation}


\label{02}


\Dot x=A_0(t)x+u_1A_1(t)x+\dotsb+u_r A_r(t)x,\quad (t,x)\in\R^{1+n},


\end{equation}


где $A_i(t)=\partial v_i(t,0)/\partial x$, равномерно локально достижимо или


равномерно согласованно (см. ниже теорему~1 и следствие~1).


Показано также, что свойства равномерной достижимости и согласованности


тесно связаны с так называемым ``методом поворотов''


В.М.\,Миллионщикова, применявшимся при исследовании поведения показателей


Ляпунова при малых возмущениях параметров уравнения.





\section{Обозначения и определения}





Пусть $\R^n$ --- евклидово пространство размерности $n$, $|x|=\sqrt{x^*x}$


--- норма в $\R^n$ ($*$ --- операция транспонирования). Если не оговорено


другое, векторы-столбцы обозначаются латинскими буквами, векторы-строки ---


греческими (таким образом, запись $\xi x$ означает скалярное произведение


векторов $\xi$ и $x$); $\cal O_{\ve}^n(x)=\{y\in\R^n:|y-x|\le\ve\}$,


$\cal O_{\ve}^n=\cal O_{\ve}^n(0)$, $S^{n-1}=\{x\in\R^n:|x|=1\}$. Для


произвольного множества $D\subset\R^n$ через $\overline D$ обозначим 


замыкание, $\Int D$ --- внутренность, $\cal O_{\ve}^n(D)$ ---


$\ve$-окрестность, $\xi\to c(\xi,D)$ --- опорную функцию


(\cite{L}, гл.\,2, \S\,12) множества $D$.





Пространство $M_{n,m}$ линейных операторов из $\R^m$ в $\R^n$ будем


отождествлять с пространством матриц размерности $n\times m$ (если $n=m$, то


пишем $M_n$); $|A|=\max\{|Ax|: |x|=1\}$ --- норма в $M_{n,m}$.  Для


$Q\in M_{n,m}$ обозначим $\cal B_{\ve}(Q)=\{H\in M_{n,m}:|H-Q|\le\ve\}$,


$\cal B_{\ve}=\cal B_{\ve}(0)$ (индексы $m$ и $n$ в обозначении


$\cal B_{\ve}(Q)$ опускаем, если ясно, о каком пространстве идет речь).





Введем топологическую динамическую систему $(\Omega,f^t)$


($\Omega$ --- полное  метрическое сепарабельное пространство, $f^t$ --- поток


на $\Omega$, т.\,е.~однопараметрическая группа преобразований $\Omega$ в себя,


непрерывная по $(t,\omega)\in\R\times\Omega$). Обозначим через $\gamma(\omega)$


траекторию движения $t\to f^t\omega$, $\gamma_+(\omega)$ --- положительную


полутраекторию, $\overline{\gamma}(\omega)$ --- замыкание $\gamma(\omega)$ в


метрике пространства $\Omega$. Напомним, что множество $E\subset\Omega$


называется инвариантным, если $f^t E\subset E$ при всех $t\in\R$, и $E$ называется


минимальным, если оно инвариантно, компактно и $\overline\gamma(\omega)=E$


для каждого $\omega\in E$. Если $E$ минимально, то для каждого $\omega\in E$


имеет место равенство $\overline\gamma(\omega)=\overline\gamma_+(\omega)$. Отметим,


что всякое компактное инвариантное множество $E$ в $\Omega$ содержит по


крайней мере одно минимальное подмножество; далее, если


$\overline{\gamma}_+(\omega)$ компактно, то омега-предельное множество


$L^+(\omega)$ точки $\omega$ непусто (при этом $L^+(\omega)$ инвариантно и


компактно), и для всякого компактного $E$ определена зона притяжения


$W^s(E)\doteq\{\omega\in\Omega:L^+(\omega)\subset E\}$ множества $E$.





Каждой паре $(\A,\omega)$, где


$\A\doteq(A_0,A_1,\dotsc,A_r):\Omega\to M_{n,m}$,


$m=n(r+1)$, и вектору $u=(u_1,\dotsc,u_r)$ поставим в соответствие


уравнение


\begin{equation}


\label{1}


\Dot x=A_0(f^t\omega)x+u_1A_1(f^t\omega)x+\dotsb+


u_rA_r(f^t\omega)x,\quad x\in\R^n.


\end{equation}


Будем предполагать, что для каждого $\omega_0\in\Omega$ функция


$t\to|\A(f^t\omega_0)|$ измерима по Лебегу, ограничена на $\R$ и для любых


$\ve>0$ и $N>0$ найдется $\delta>0$ такое, что выполнено неравенство


$\max\limits_{|t|\le N}\int\limits_t^{t+1}


|\A(f^s\omega)-\A(f^s\omega_0)|ds<\ve$,


как только $\rho(\omega,\omega_0)<\delta$ ($\rho$ --- метрика в $\Omega$).


Уравнение (\ref{1}) будем отождествлять с парой $(\A,\omega)$.





Пусть задано произвольное множество $U\subset\R^r$ такое, что $0\in U$;


множество $\cal U$, состоящее из всех измеримых функций $u:\R\to U$, будем


называть множеством допустимых управлений. Зафиксируем уравнение


$(\A,\omega)$.


Пусть $X_u(t,s,\omega)$ --- функция Коши уравнения (\ref{1}) при


фиксированном управлении $u=u(t,\omega)$.  Отметим, что функция


$\omega\to X_u(t,s,\omega)$ непрерывна в каждой точке $\omega_0\in\Omega$


равномерно относительно $(t,s)$ на любом компакте в $\R^2$ и, если


управление стационарно относительно потока,


т.\,е.~$u(t,\omega)=u(f^t\omega)$,


то $X_u(t+\tau,s+\tau,\omega)=X_u(t,s,f^{\tau}\omega)$.





Для любого $\vt>0$ построим множество достижимости


$$


\frak D_{\vt}(\omega)=\{X\in M_n: X=X_u(\vt,0,\omega),


\ u(\cdot)\in\cal U\}


$$


уравнения (\ref{1}) в пространстве матриц $M_n$ из единичной матрицы и


множество


$$


\cal B_{\ve}(I)X_0(\vt,0,\omega)=\{X\in M_n: X=HX_0(\vt,0,\omega),\


H\in\cal B_{\ve}(I)\},


$$


где $I$ --- единичная матрица.





\begin{lem}


\label{l1}


Имеют место включения


\begin{equation}


\label{2}


\cal B_{\nu_1}(X_0(\vt,0,\omega))\subset


\cal B_{\ve}(I)X_0(\vt,0,\omega)\subset


\cal B_{\nu_2}(X_0(\vt,0,\omega)),


\end{equation}


где $\nu_1=\ve/|X_0(0,\vt,\omega)|$, $\nu_2=\ve|X_0(\vt,0,\omega)|$.


\end{lem}





\begin{pf} Пусть $H=I+G$, тогда


$$


\cal B_{\ve}(I)X_0(\vt,0,\omega)=


X_0(\vt,0,\omega)+\cal B_{\ve}X_0(\vt,0,\omega),


$$


где


$\cal B_{\ve}X_0(\vt,0,\omega)=


\{X\in M_n: X=GX_0(\vt,0,\omega), \ |G|\le\ve\}$.


Поэтому включения (\ref{2}) эквивалентны включениям


$\cal B_{\nu_1}\subset\cal B_{\ve}X_0(\vt,0,\omega)\subset


\cal B_{\nu_2}$. Если


$A\in\cal B_{\ve}$, то $|AX_0(\vt,0,\omega)|\le\ve|X_0(\vt,0,\omega)|=\nu_2$,


откуда следует второе включение. Если $|A|\le\nu_1$, то


$|AX_0(0,\vt,\omega)|\le\ve$, поэтому


$$


A=AX_0(0,\vt,\omega)X_0(\vt,0,\omega)\in\cal B_{\ve}(I)X_0(\vt,0,\omega),


$$


откуда следует первое включение.


\end{pf}





\begin{defn}


\label{d1}


Уравнение (\ref{1}) называется


\begin{itemize}


\item[а)] локально достижимым, если найдутся $\vt>0$ и $\ve>0$ такие, что


выполнено включение


$$


\cal B_{\ve}(I)X_0(\vt,0,\omega)\subset\frak D_{\vt}(\omega);


$$


\item[б)] равномерно локально достижимым, если найдутся $\vt>0$ и $\ve>0$


такие, что для всех $\omega_0\in\ov{\gamma}(\omega)$


имеет место включение


$$


\cal B_{\ve}(I)X_0(\vt,0,\omega_0)\subset\frak D_{\vt}(\omega_0).


$$


\end{itemize}


\end{defn}





\begin{lem}


\label{l2}


Уравнение $(\ref{1})$ локально достижимо в том и только том


случае, если существуют такие $\vt>0$ и $\ve>0$, что найдется


управление $(t,H)\to u(t,H)$, определенное на $[0,\vt]\times\cal B_{\ve}(I)$


и обеспечивающее равенство


\begin{equation}


\label{H}


X_u(\vt,0,\omega)=HX_0(\vt,0,\omega).


\end{equation}


Уравнение $(\ref{1})$ равномерно локально достижимо в том и


только том случае, если существуют такие $\vt>0$ и $\ve>0$, что для


любой непрерывной функции $H:\ov{\gamma}(\omega)\to\cal B_{\ve}(I)$


найдется управление $(t,\omega_0)\to u(t,\omega_0)$, определенное на


$[0,\vt]\times\ov{\gamma}(\omega)$ и обеспечивающее равенство


\begin{equation}


\label{Hom}


X_u(\vt,0,\omega_0)=H(\omega_0)X_0(\vt,0,\omega_0).


\end{equation}


\end{lem}





\begin{pf} Пусть при $u=u(t,H)$ выполнено равенство


(\ref{H}); тогда для всех $H\in\cal B_{\ve}(I)$ имеем


$HX_0(\vt,0,\omega)\in\frak D_{\vt}(\omega)$, следовательно,


уравнение (\ref{1}) локально достижимо.





Очевидно, что из равномерной локальной достижимости следует


(\ref{Hom}). Пусть для всякой непрерывной $H(\omega_0)$ выполнено


равенство (\ref{Hom}). Если найдется $H_0\in\cal B_{\ve}(I)$ такая, что


матрица $H_0X_0(\vt,0,\omega_0)$ не содержится в


$\frak D_{\vt}(\omega_0)$, то для функции $H(\omega_0)\equiv H_0$ и любого


допустимого управления равенство (\ref{Hom}) не выполнено.


\end{pf}





\begin{note}


\label{rem1}


Наличие свойства локальной или равномерной локальной достижимости


предоставляет возможность (в силу (\ref{H}) или (\ref{Hom}))


``немного поворачивать'' функцию Коши уравнения (\ref{1}) с помощью


подходящего управления (и тем самым влиять на поведение решений). Не


всякое уравнение обладает этим свойством, но очевидно, что уравнение


$\Dot x=A(t)x+U(t)x$, где матрица $A(t)$ ограничена, а элементы матрицы


$U(t)$ интерпретируются как управляющие функции ($|U(t)|\le\ve$),


равномерно локально достижимо. Таким образом, выполнено  (\ref{Hom}) и


это обстоятельство позволило В.М.\,Миллионщикову \cite{M69S}, \cite{M69D}, а


вслед за ним и другим исследователям \cite{Izob}, построить


современную теорию показателей Ляпунова. Сам же метод


``поворотов'' получил название метода поворотов Миллионщикова.


\end{note}





По уравнению $(\A,\omega)$ построим матричное дифференциальное уравнение


\begin{equation}


\label{sys3}


\Dot Z=A_0(f^t\omega)Z+


\big(u_1A_1(f^t\omega)+\dotsb+u_rA_r(f^t\omega)\big)X_0(t,0,\omega)


\end{equation}


относительно $Z\in M_n$. Обозначим через $\frak L_{\vt}(\omega)$ множество


достижимости уравнения (\ref{sys3}) из нуля за время $\vt$ под


действием управлений $u:[0,\vt]\to\R^r$. Таким образом,


$G\in\frak L_{\vt}(\omega)$ в том и только том случае, если найдется


$u_G:[0,\vt]\to\R^r$, что уравнение (\ref{sys3}) при $u=u_G(t)$


имеет решение $Z(t)$, удовлетворяющее условиям $Z(0)=0$, $Z(\vt)=G$. В


силу линейности уравнения (\ref{sys3}) множество $\frak L_{\vt}(\omega)$


является линейным подпространством в $M_n$.





\begin{defn}


\label{d2}


Уравнение $(\A,\omega)$ называется согласованным, если найдется $\vt>0$


такое, что $\frak L_{\vt}(\omega)=M_n$, и равномерно согласованным, если


найдутся $\vt>0$ и $\beta>0$ такие, что для всех


$\omega_0\in\ov{\gamma}(\omega)$ множество


$\frak L_{\vt}(\omega_0)$ совпадает с


$M_n$, и выполнено неравенство


$|u_G(t,\omega_0)|\le\beta|G|$, $0\le t\le\vt$


(здесь $u_G(t,\omega_0)$ --- управление, переводящее решение уравнения


(\ref{sys3}) при $\omega=\omega_0$ из нуля в точку $G$ в момент $t=\vt$).


\end{defn}





\begin{note}


\label{rem2}


Непосредственно из определения следует, что если уравнение $(\A,\omega)$


согласованно (равномерно согласованно), то для любого $\tau\in\R_+$


($\tau\in\R$) уравнение $(\A,f^{-\tau}\omega)$ согласованно (равномерно


согласованно).


\end{note}





\begin{note}


\label{rem3}


Определения согласованности и равномерной согласованности введены в


\cite{PT94} для системы


$\Dot x=(A(f^t\omega)+B(f^t\omega)UC^*(f^t\omega))x$


(элементы матрицы $U$ интерпретируются как управляющие функции)


в связи с задачами управления показателями Ляпунова и изучались в


\cite{PT95}--\cite{PT97} (см. библиографию в \cite{PT97}).


\end{note}





\section{Равномерная локальная достижимость}





Основное свойство равномерно локально достижимых уравнений содержится в


формулируемой ниже теореме \ref{th1}. Эта теорема утверждает, что при


надлежащем выборе допустимого управления все решения уравнения (\ref{1})


ведут себя подобно решениям заранее заданного модельного уравнения с


дискретным временем. Таким образом, появляется возможность влиять на


асимптотические свойства решений уравнения (\ref{1}), например,


управлять показателями Ляпунова или аналогичными характеристиками,


отвечающими за асимптотическое поведение решений. В прикладных задачах


это бывает важно, потому что позволяет разделять движения (добиваясь


структурной устойчивости) и улучшать асимптотику движений за счет


перемещения показателей влево. Если же модельное уравнение стационарно:


$y_{k+1}=Qy_k$ (см. (\ref{Q}) в формулируемой ниже теореме~\ref{th1}),


то уравнение (\ref{1}) ``почти стационарно'', т.\,е.~приводимо


ляпуновским преобразованием к стационарному уравнению (следствие


\ref{St}).





Предварительно докажем две простые леммы.





\begin{lem}


\label{lm3}


Пусть $E$ --- инвариантное компактное множество в $\Omega$ и для


каждого $\omega\in E$  уравнение $(\ref{1})$ равномерно локально


достижимо. Тогда оно равномерно локально достижимо на множестве $E$,


т.\,е.~найдутся общие для всех $\omega\in E$ константы $\vt>0$ и $\ve>0$,


участвующие в определении равномерной локальной достижимости.


\end{lem}





\begin{pf} Пусть $\vt(\omega)$ и $\ve(\omega)$ --- константы,


участвующие в определении равномерной локальной достижимости уравнения


$(\A,\omega)$. Надо показать, что $\sup_E\theta(\omega)<\fy$, где


$\theta(\omega)=\max\{\vt(\omega),\ve^{-1}(\omega)\}$.


Если $\sup_E\theta(\omega)=\fy$,


то найдется такая сходящаяся (в силу компактности $E$) последовательность


$\{\omega_i\}$, что $\theta(\omega_i)\to\fy$. Это означает, что уравнение


$(\A,\wh\omega)$, где $\wh\omega=\lim\omega_i$,


не является равномерно локально достижимым.


\end{pf}





\begin{lem}


\label{lmS}


Пусть $E$ --- инвариантное компактное множество в $\Omega$,


$u:\R\times E\to U$ --- произвольное допустимое управление. Если для


некоторого $\tau>0$ и всех $t\in\R_+$ имеет место равенство


\begin{equation}


\label{S}


u(t,\omega)=u(t-\tau,f^{\tau}\omega),


\end{equation}


то $X_u(t+\tau,s+\tau,\omega)=X_u(t,s,f^{\tau}\omega)$,


$(t,s)\in\R_+\times\R_+$.


Если равенство $(\ref{S})$ выполнено для любого $\tau\in\R$ и всех


$t\in\R$, то управление $u(t,\omega)$ стационарно относительно потока


$f^t$, т.\,е.~$u(t,\omega)=\wh u(f^t\omega)$ для некоторой функции


$\wh u:E\to U$.


\end{lem}





\begin{pf} Введем в рассмотрение матричное уравнение


\begin{equation}


\label{X}


\Dot X=A_0(f^t\omega)X+u_1A_1(f^t\omega)X+\dotsb


+u_r A_r(f^t\omega)X,\quad X\in M_n,


\end{equation}


отвечающее уравнению (\ref{1}).


Для доказательства первого утверждения леммы достаточно подставить


управление, удовлетворяющее равенству (\ref{S}), в уравнение


(\ref{X}) и провести замену времени $t\to t+\tau$. Тогда в силу


определения $u(t,\omega)$ получим равенство


$u(t+\tau,\omega)=u(t,f^{\tau}\omega)$, которое эквивалентно (\ref{S}).





Второе утверждение следует из (\ref{S}) при $\tau=t$.


\end{pf}





\begin{thm}


\label{th1}


Пусть $E$ --- инвариантное компактное множество в $\Omega$ и для


каждого $\omega\in E$  уравнение $(\ref{1})$ равномерно локально


достижимо. Тогда существуют $\vt>0$ и $\ve>0$ такие, что для любой


функции $Q:E\to M_n$, удовлетворяющей при всех $\omega\in E$ неравенству


$|Q(\omega)X_0(0,\vt,\omega)-I|\le\ve$, найдется управление


$u:\R\times E\to U$,


обеспечивающее для решения $x_u(t,x_0,\omega)$ уравнения $(\ref{1})$


при $u=u(t,\omega)$, всех $x_0\in\R^n$ и всех $k=0,1,\dotsc$ равенства


$x_u(t_k,x_0,\omega)=y_k(x_0,\omega)$, где $t_k=k\vt$,


$\{y_k(x_0,\omega)\}_{k=0}^{\fy}$ --- решение задачи


\begin{equation}


\label{Q}


y_{k+1}=Q(f^{t_k}\omega)y_k,\quad y_0=x_0,\quad k=0,1,\dotsc


\end{equation}


\end{thm}





{\bf Доказательство.} Пусть $\vt$ и $\ve$ как в лемме \ref{lm3}, и


$Q(\omega)$ удовлетворяет условию теоремы. Для каждого $\omega\in E$ построим


управление $u^0(\cdot,\omega):[0,\vt]\to U$, определенное на отрезке


$[0,\vt]$ и такое, что (см. лемму \ref{l2}) при $u=u^0(t,\omega)$ для


уравнения (\ref{1}) выполнено равенство


$X_{u^0}(\vt,0,\omega)=H(\omega)X_0(\vt,0,\omega)$, где


$H(\omega)=Q(\omega)X_0(0,\vt,\omega)$. Следовательно,


для любого $x_0\in\R^n$


имеет место равенство $x_{u^0}(\vt,x_0,\omega)=Q(\omega)x_0$ или в силу


(\ref{Q}) $x_{u^0}(\vt,x_0,\omega)=y_1(x_0,\omega)$.





Тем самым для каждого $k=0,1,\dotsc$ построено управление


$u^k=u^0(t,f^{t_k}\omega)$, определенное на $[0,\vt]$ и обеспечивающее


для решений уравнения $(\A,f^{t_k}\omega)$ равенство


$x_{u^k}(\vt,x_0,f^{t_k}\omega)=y_1(x_0,f^{t_k}\omega)$.





``Растянем'' последовательность $\{u^k(t)\}^{\fy}_{k=0}$,


$t\in [0,\vt]$, на $\R_+$, определив управление $u:\R_+\times E$ на


каждом полуинтервале $\Delta_k=[t_k,t_{k+1})$ равенством


$u(t,\omega)=u^0(t-t_k,f^{t_k}\omega)$ при $t\in\Delta_k$. Непосредственной


проверкой можно убедиться, что для каждого целого $m\ge 0$ при


$\tau=t_m$ и всех $t\in\R_+$ выполнено равенство (\ref{S}). Поэтому


в силу леммы \ref{lmS} имеют место равенства


$X_u(t_{k+1},t_k,\omega)=X_u(\vt,0,f^{t_k}\omega)$.





Пусть $x_u(t,x_0,\omega)$ --- решение уравнения $(\A,\omega)$, отвечающее


управлению $u(t,\omega)$. Покажем тогда, что последовательность


$\{x_u(t_k,x_0,\omega)\}_{k=1}^{\fy}$ является решением задачи (\ref{Q}).


Действительно,


$$


x_u(t_{k+1},x_0,\omega)=X_u(t_{k+1},0,\omega)x_0=


X_u(t_{k+1},t_k,\omega)X_u(t_k,0,\omega)x_0.


$$


С учетом равенства $X_u(t_{k+1},t_k,\omega)=X_u(\vt,0,f^{t_k}\omega)$ имеем


далее


$$


x_u(t_{k+1},x_0,\omega)=X_u(\vt,0,f^{t_k}\omega)x_u(t_k,x_0,\omega)=


Q(f^{t_k}\omega)x_u(t_k,x_0,\omega).\quad\square


$$





\begin{cor}


\label{St}


Пусть множество $\ov\gamma(\omega)$ компактно и


уравнение (\ref{1}) равномерно локально достижимо. Тогда существуют


такие $\vt>0$ и $\ve>0$, что для любой вещественной матрицы


$\varLambda\in M_n$, удовлетворяющей неравенству


$|\exp(\vt\varLambda)X_0(0,\vt,\omega)-I|\le\ve$, найдется управление


$t\to u(t,\omega)\in U$, при котором


уравнение (\ref{1}) приводимо ляпуновским преобразованием к


уравнению $\Dot y=\varLambda y$ с постоянной матрицей $\varLambda$.





Управление $t\to u(t,\omega)$ обладает следующим свойством:


если $\omega$ --- периодическая точка (т.\,е.~найдется $T>0$, для


которого $f^{t+T}\omega\equiv f^t\omega$), то при некотором целом $m\ge 1$


управление периодично с периодом $mT$: $u(t+mT,\omega)\equiv u(t,\omega)$.


\end{cor}





\begin{pf} Пусть $Q(\omega)\equiv\exp(\vt\varLambda)$, тогда


матрица $Q$ удовлетворяет условиям теоремы \ref{th1}. По матрице $Q$


построим управление $t\to u(t,\omega)$ (как это было сделано в


доказательстве теоремы~\ref{th1}) и матрицу


$\varPhi(t)=X_u(t,0,\omega)\exp(-t\varLambda)$. В силу известной теоремы


Н.П.\,Еругина (\cite{Er}, \S\,19, с.\,119) достаточно показать, что


матрица $\varPhi(t)$ является ляпуновской матрицей.





Функция $t\to\varPhi(t)$ непрерывна и ограничена на $\R_+$. Действительно,


из (\ref{Q}) для всех $x_0\in\R^n$ имеем равенство


$x_u(t_k,x_0,\omega)=\exp(t_k\varLambda)x_0$, поэтому


$X_u(t_k,0,\omega)=\exp(t_k\varLambda)$. В силу компактности


$\ov\gamma(\omega)$, ограниченности $\A(f^t\omega)$ и управления


$u(t,\omega)$ на $\R_+$ найдется такая константа $a$, что


$|X_u(t,s,f^{\tau}\omega)|\le a$ для всех


$(t,s)\in\Delta_k\times\Delta_k$, всех


$k=0,1,\dotsc$ и всех $\tau\in\R_+$. Пусть $t\in\Delta_k$, тогда


$\varPhi(t)=X_u(t,t_k,\omega)X_u(t_k,0,\omega)\exp(-t\varLambda)=


X_u(t,t_k,\omega)\exp(t_k-t)\varLambda$,


поэтому $|\varPhi(t)|\le c$ для всех


$t\in\R_+$, причем $c$ не зависит от $k$. Аналогично доказывается


ограниченность $\varPhi^{-1}(t)$.





Ограниченность производной $\Dot\varPhi(t)$ следует из равенства


$$


\Dot\varPhi(t)=(A_0(f^t\omega)+u_1(t,\omega)A_1(f^t\omega)+\dotsb+


u_r(t,\omega)A_r(f^t\omega))\varPhi(t)-\varPhi(t)\Lambda


$$


и ограниченности $\varPhi(t)$.





Докажем второе утверждение следствия. Отметим, во-первых, что


непосредственно из определения равномерной локальной достижимости и


условия $0\in U$ следует, что если это свойство имеет место при


некоторых $\vt$ и $\ve$, то оно имеет место при всех $\vt_1\ge\vt$.


Поэтому найдется такое целое $m\ge 1$, что уравнение~(\ref{1}) равномерно


локально достижимо при $\vt=mT$. В силу (\ref{S}) имеем


$u(t+mT,\omega)=u(t,f^{mT}\omega)=u(t,\omega)$.


\end{pf}





\section{Равномерная согласованность и достижимость}





В этом параграфе показано, что в некритическом случае ($0\in\Int U$)


равномерная согласованность --- это равномерная локальная достижимость


по первому приближению.





В силу леммы \ref{l2} локальная достижимость уравнения (\ref{1})


эквивалентна разрешимости уравнения (\ref{X}) с условием (\ref{H}), а


равномерная локальная достижимость --- с условием (\ref{Hom}).


Произведем в (\ref{X}) замену $Y(t,\omega)=X_u(t,0,\omega)-X_0(t,0,\omega)$,


тогда уравнение относительно $Y$ запишется в виде


\begin{multline}


\label{Y}


\Dot Y=A_0(f^t\omega)Y+\Big(u_1A_1(f^t\omega)+\dotsb


+u_rA_r(f^t\omega)\Big)Y+\\


+\Big(u_1A_1(f^t\omega)+\dotsb+u_rA_r(f^t\omega)\Big)


X_0(t,0,\omega),\quad Y\in M_n,


\end{multline}


а условие (\ref{H}) разбивается на два условия:


$Y(0,\omega)=0$, $Y(\vt,\omega)=GX_0(\vt,0,\omega)$, где $G=H-I$. Обозначим


через $\Bbb D_{\vt}(\omega)$ множество достижимости за время $\vt>0$


уравнения (\ref{Y}) из точки $Y=0$ под действием управлений $u\in\cal U$


(отметим, что $0\in\Bbb D_{\vt}(\omega)\subset M_n$). Тогда включение


$0\in\Int\Bbb D_{\vt}(\omega)$ при некотором $\vt>0$ эквивалентно


локальной достижимости уравнения (\ref{1}), а включение


$\cal B_{\ve}\subset\Bbb D_{\vt}(\omega_0)$, выполненное при некоторых


$\vt>0$,


$\ve>0$ и всех $\omega_0\in\ov{\gamma}(\omega)$, эквивалентно равномерной


локальной достижимости.





\begin{thm}


\label{th2}


Пусть $0\in\Int U$ и уравнение $(\ref{1})$ согласованно


$($равномерно согласованно$)$, тогда оно локально достижимо


$($равномерно локально достижимо$)$.


\end{thm}





\begin{pf} Уравнение (\ref{sys3}) служит линейным


приближением (в окрестности точки $u=0$, $Y=0$) для уравнения (\ref{Y}),


а свойство согласованности, сформулированное в терминах уравнения


(\ref{sys3}), эквивалентно свойству полной достижимости уравнения


(\ref{sys3}) (т.\,е.~множество достижимости $\frak L_{\vt}(\omega)$ уравнения


(\ref{sys3}) совпадает с пространством $M_n$). Воспользуемся теоремой 3


работы \cite{T74}. С этой целью матрицы $Y$ и $Z$ в уравнениях (\ref{Y})


и (\ref{sys3}) ``развернем в вектор'', а в теореме 3 из


\cite{T74} проведем замену времени $t\to\vt-t$ и слова ``локальная


управляемость в малом'' заменим на слова ``локальная достижимость


в малом'' (локальная достижимость в малом означает локальную


достижимость ``развернутого уравнения'' (\ref{Y}), дополненную


следующим свойством: для любого $\ve>0$ найдется $\delta>0$ такое, что


если $|G|<\delta$, то решение $Y(t,\omega)$, переходящее в точку


$Y(\vt,\omega)=G$


при надлежащем управлении, удовлетворяет для всех $t\in[0,\vt]$


неравенству $|Y(t,\omega)|<\ve$). В силу условия $0\in\Int U$ и цитируемой


теоремы 3 из \cite{T74} первое утверждение доказано.





Второе утверждение теоремы после аналогичных рассуждений следует из


теоремы 1 из \cite{T82} (в этой работе требуется минимальность


множества $\ov\gamma(\omega)$, но это предположение излишне:


достаточно предполагать, что $\ov\gamma(\omega)$ компактно).


\end{pf}





\section{Критерии согласованности}





В этом параграфе результаты работ \cite{PT94}, \cite{PT97} о согласованности


уравнения


\begin{equation}


\label{4.0}


\Dot x=\big(A(f^t\omega)+B(f^t\omega)UC^*(f^t\omega)\big)x,\quad x\in\R^n,


\end{equation}


где $U$ --- $(m\times k)$-матрица управляющих параметров,


переносятся на уравнение (\ref{1}).


Для каждых $i,j,p,s\in\{1,\dotsc,n\}$ введем обозначения


$$


\gamma_{ijps}(\vt,\omega)=\int_0^{\vt}\sum_{l=1}^re_i^*\wh


A_l(t,\omega)e_p e_s^*\wh A_l^*(t,\omega)e_j\,dt,


$$


где $e_k\in\R^n$ --- $k$-й столбец единичной матрицы,


$\wh A_l(t,\omega)=X_0(0,t,\omega)A_l(f^t\omega)X_0(t,0,\omega)$,


$l=1,\dotsc,r$. Построим далее $(n\times n)$-матрицы


$\Gamma_{ij}(\vt,\omega)=\{\gamma_{ijps}(\vt,\omega)\}_{p,s=1}^n$,


$i,j=1,\dotsc,n$, и $(n^2\times n^2)$-матрицу


$\Gamma(\vt,\omega)=\{\Gamma_{ij}(\vt,\omega)\}_{i,j=1}^n$,


которую в дальнейшем будем называть матрицей согласования (уравнения


(\ref{1}) на $[0,\vt]$).





\begin{lem}


\label{l4.1}


Матрица согласования обладает следующими свойствами$:$





{\em а)} $\Gamma(\vt,\omega)=\Gamma^*(\vt,\omega);$





{\em б)} $\lambda(\vt,\omega)\ge 0$, $(\lambda$ --- наименьшее собственное значение


матрицы $\Gamma);$





{\em в)} $\lambda(\vt_1,\omega)\ge\lambda(\vt,\omega)$ при $\vt_1\ge\vt$.


\end{lem}





\begin{pf} Утверждение а) следует из равенства


$\gamma_{ijps}=\gamma_{jisp}$. Докажем утверждение б). Достаточно доказать,


что для любого $h=\col(h_{11},\dotsc,h_{1n},\dotsc,h_{n1},\dotsc,h_{nn})$


справедливо неравенство $h^*\Gamma h\ge 0$. Имеем


\begin{multline}


\label{4.1}


h^*\Gamma h=\sum_{i,j,p,s=1}^n h_{ip}h_{js}\gamma_{ijps}=\int_0^{\vt}


\sum_{l=1}^r\;\sum_{i,j,p,s=1}^n h_{ip}h_{js}e_i^*\wh A_l(t,\omega)e_pe_s^*


\wh A_l^*(t,\omega)e_j\,dt=\\


=\int_0^{\vt}\sum_{l=1}^r


\bigg(\,\sum_{i,p=1}^n h_{ip}e_i^*\wh A_l(t,\omega)e_p\bigg)


\bigg(\,\sum_{j,s=1}^nh_{js}e_s^*\wh A_l^*(t,\omega)e_j\bigg)\,dt=\\


=\int_0^{\vt}\sum_{l=1}^r\bigg(\,\sum_{i, p=1}^n


h_{ip}e_i^*\wh A_l(t,\omega)e_p\bigg)^2\,dt\ge 0.


\end{multline}


Утверждение в) следует из неравенства (\ref{4.1}).


\end{pf}





Определим функции $u_{ip}^l:[0,\vt]\times\Omega\to\R$ и вектор-функции


$u_{ip}$ равенствами


\begin{gather}


\label{4.2}


u_{ip}^l(t,\omega)=e_i^*\wh A_l(t,\omega)e_p,


\quad i,p=1,\dotsc,n,\quad l=1,\dotsc,r,\\


\label{4.3}


u_{ip}(t,\omega)=\col(u_{ip}^1(t,\omega),\dotsc,u_{ip}^r(t,\omega)),


\quad i,p=1,\dotsc,n.


\end{gather}





\begin{lem}


\label{l4.2}


Совокупность вектор-функций $(\ref{4.3})$ линейно независима на


$[0,\vt]$ в том и только том случае, если $\lambda(\vt,\omega)>0$.


\end{lem}





\begin{pf} Пусть функции (\ref{4.3}) линейно независимы на


$[0,\vt]$, но $\lambda(\vt,\omega)=0$. Тогда существует такой ненулевой вектор


$h$, что $\Gamma(\vt,\omega)h=0$. Следовательно, $h^*\Gamma(\vt,\omega)h=0$ и в силу


(\ref{4.1}) для всех $l=1,\dotsc,r$ выполнены равенства


\begin{equation}


\label{4.4}


\sum_{i,p=1}^n h_{ip}e_i^*\wh A_l(t,\omega)e_p=0,\quad t\in[0,\vt].


\end{equation}


Это эквивалентно равенству $\sum\limits_{i, p=1}^nh_{ip}u_{ip}(t,\omega)=0$,


$t\in[0,\vt]$, что противоречит линейной независимости функций


(\ref{4.3}).





Докажем теперь достаточность условий леммы. Если функции (\ref{4.3})


линейно зависимы, то найдется такой ненулевой вектор $h\in\R^{n^2}$,


что $\sum\limits_{i,p=1}^n h_{ip}u_{ip}(t,\omega)=0$, $t\in[0,\vt]$.


Следовательно, для всех $t\in[0,\vt]$ и $l=1,\dotsc,r$ выполнено


равенство (\ref{4.4}), поэтому в силу (\ref{4.1}) имеем


$h^*\Gamma(\vt,\omega)h=0$, т.\,е.~$\lambda(\vt,\omega)=0$.


\end{pf}





\begin{note}


\label{rem4.1}


Матрица $\Gamma(\vt,\omega)$ представляет собой матрицу Грама для


совокупности вектор-функций $\{u_{ip}(\cdot,\omega)\}_{i,p=1}^n$ на


отрезке $[0,\vt]$.


\end{note}





\begin{thm}


\label{th4.1}


Следующие утверждения эквивалентны$:$





{\em а)} уравнение $(\ref{1})$ согласованно$;$





{\em б)} найдется $\vt>0$, при котором матрица $\Gamma(\vt,\omega)$


положительно определена $(\lambda(\vt,\omega)>0);$





{\em в)} совокупность вектор-функций $\{u_{ip}(\cdot,\omega)\}_{i,p=1}^n$,


определенных равенствами $(\ref{4.2})$, $(\ref{4.3})$, линейно


независима на $[0,\vt]$ для некоторого $\vt>0$.


\end{thm}





\begin{pf} Условия б) и в) эквивалентны в силу леммы~\ref{l4.2}.


Покажем, что из б) следует а). Пусть $\lambda(\vt,\omega)>0$ при некотором


$\vt>0$. Для любой $G\in M_n$ требуется найти управление


$u=\col(u_1,\dotsc,u_r):[0,\vt]\to\R^r$, при котором уравнение


(\ref{sys3}) имеет решение, удовлетворяющее условиям $Z(0)=0$,


$Z(\vt)=G$. Таким образом, задача о построении $u(\cdot)$ сводится к


задаче о разрешимости уравнения


\begin{equation}


\label{4.5}


\int_0^{\vt}\sum_{l=1}^r\wh A_l(t,\omega)u_l(t)\,dt=Q(\vt,\omega),


\end{equation}


где $Q(\vt,\omega)=X_0(0,\vt,\omega)G$. Решение


$\wh u=\col(\wh u_1,\dotsc,\wh u_r)$ уравнения (\ref{4.5}) будем искать


в виде $\wh u(t)=\sum\limits_{j,s=1}^n h_{js}u_{js}(t)$, где $u_{js}(t)$


определены равенствами (\ref{4.2}), (\ref{4.3}). Подставляя $\wh u(t)$


в (\ref{4.5}) и умножая (\ref{4.5}) слева на $e_i^*$, а справа на $e_p$,


получим систему $n^2$ алгебраических уравнений


\begin{equation}


\label{4.6}


\sum_{j, s=1}^n\gamma_{ijps}h_{js}=q_{ip},\quad i,p=1,\dotsc,n,


\end{equation}


относительно $h_{js}$. Здесь $q_{ip}=e_i^*Q(\vt,\omega)e_p$. Поскольку


$\lambda(\vt,\omega)>0$, то система (\ref{4.6}) разрешима:


$h=\Gamma^{-1}q$, где


$h=\col(h_{11},\dotsc,h_{1n},\dotsc,h_{n1},\dotsc,h_{nn})$,


$q=\col(q_{11},\dotsc,q_{1n},\dotsc,q_{n1},\dotsc,q_{nn})$.


Поэтому уравнение (\ref{1}) согласованно. Заметим, что для решения


$\wh u(t)$ уравнения (\ref{4.5}) имеется оценка


$$


|\wh u(t)|\le\sum_{j,s=1}^n|h_{js}|\,


|u_{js}(t)|\le|\Gamma^{-1}(\vt,\omega)|\sum_{j,s=1}^n


|u_{js}(t)|\,|q(\vt,\omega)|.


$$


Далее, из ограниченности $A_l(f^t\omega)$ и $X_0(t,s,\omega)$ при


$(t,s)\in[0,\vt]\times[0,\vt]$ следует ограниченность на $[0,\vt]$


функций $u_{js}(t)$. Поэтому найдется такая константа $\beta$ (не зависящая от


$G$, но зависящая от $\vt$), что $|\wh u(t)|\le\beta|G|$, $0\le t\le\vt$.





Докажем, что из а) следует в). Выберем $\vt>0$ из определения


согласованности. Достаточно показать, что если уравнение (\ref{4.5})


разрешимо относительно $u(t)$ для любой $Q$, то совокупность


вектор-функций (\ref{4.3}) линейно независима на $[0,\vt]$. Предположим


противное, пусть существует ненулевой вектор


$h=\col(h_{11},\dotsc,h_{1n},\dotsc,h_{n1},\dotsc,h_{nn})$ такой, что


$\sum\limits_{i, p=1}^n h_{ip}u_{ip}(t,\omega)=0$, $t\in[0,\vt]$. Тогда


$\sum\limits_{i,p=1}^n h_{ip}e_i^*\wh A_l(t,\omega)e_p=0$ при всех


$l=1,\dotsc,r$ и $t\in[0,\vt]$. Для матрицы $H=\{h_{ip}\}_{i,p=1}^n$


существует такая функция $\wh u=\col(\wh u_1,\dotsc,\wh u_r):[0,\vt]\to\R^r$,


что


$$


\int_0^{\vt}X_0(0,t,\omega)\big(\wh u_1(t)A_1(f^t\omega)+\dotsb+


\wh u_r(t)A_r(f^t\omega)\big)X_0(t,0,\omega)\,dt=H,


$$


поэтому для всех $i, p=1,\dotsc,n$


$$


\int_0^{\vt}e_i^*X_0(0,t,\omega)\big(\wh u_1(t)A_1(f^t\omega)+\dotsb+


\wh u_r(t)A_r(f^t\omega)\big)X_0(t,0,\omega)e_p\,dt=h_{ip}.


$$


Умножая каждое из этих равенств на $h_{ip}$ и суммируя по


$i,p=1,\dotsc,n$, получим


\begin{gather*}


\sum_{i,p=1}^nh_{ip}^2=\sum_{i,p=1}^nh_{ip}\,


\int_0^{\vt}e_i^*X_0(0,t,\omega)\big(\wh u_1(t)A_1(f^t\omega)+\dotsb+


\wh u_r(t)A_r(f^t\omega)\big)X_0(t,0,\omega)e_p\,dt=\\


=\int_0^{\vt}\,\sum_{l=1}^r\wh u_l(t)


\bigg(\,\sum_{i, p=1}^n h_{ip}e_i^*\wh A_l(t,\omega)e_p\bigg)\,dt=0,


\end{gather*}


что противоречит условию $h\ne 0$.


\end{pf}





\begin{note}


\label{rem4.2}


Пусть задана тройка $(A_0(\omega),B(\omega),C(\omega))$, где $B=(b_1,\dotsc,b_m)$,


$b_i\in\R^n$, $C=(c_1,\dotsc,c_k)$, $c_j\in\R^n$. Положим $r=mk$ и


построим матрицы $A_1=b_1c_1^*$, $A_2=b_1c_2^*,\dotsc,\,A_k=b_1c_k^*$,


$A_{k+1}=b_2c_1^*,\dotsc,\,A_{2k}=b_2c_k^*,\dotsc,\,A_{r-k+1}=b_mc_1^*,


\dotsc,\,A_r=b_mc_k^*$. Тогда определение согласованности для уравнения


$\Dot x=A_0(f^t\omega)x+u_1A_1(f^t\omega)x+\dotsb+u_r A_r(f^t\omega)x$ с


построенными таким образом матрицами $A_l$ совпадает с определением


согласованности, введенным в \cite{PT94} для уравнения (\ref{4.0}).


\end{note}





\begin{thm}


\label{th4.2}


Пусть множество $\overline{\gamma}(\omega)$ компактно. Тогда уравнение


$(\A,\omega)$ равномерно согласованно в том и только том случае, если


найдутся $\vt>0$ и $\ve>0$ такие, что для всех $\omega_0\in\overline\gamma(\omega)$


наименьшее собственное значение $\lambda(\vt,\omega_0)$ матрицы согласования


$\Gamma(\vt,\omega_0)$ удовлетворяет неравенству $\lambda(\vt,\omega_0)\ge\ve$.


\end{thm}





\begin{pf} Пусть уравнение $(\A,\omega)$ равномерно


согласованно. Поскольку множество $\overline{\gamma}(\omega)$ компактно, а


функция $\omega_0\to\lambda(\vt,\omega_0)$ непрерывна, то существует


$\omega_1\in\overline{\gamma}(\omega)$ такое, что


$\lambda(\vt,\omega_1)=\min\{\lambda(\vt,\omega_0):


\omega_0\in\overline{\gamma}(\omega)\}$.


Если $\lambda(\vt,\omega_1)=0$, то уравнение $(\A,\omega_1)$ не является


согласованным, что противоречит определению \ref{d2}.





Пусть $\lambda(\vt,\omega_0)\ge\ve$ при некоторых $\vt>0$, $\ve>0$ и всех


$\omega_0\in\overline\gamma(\omega)$. Тогда для каждого


$\omega_0\in\overline\gamma(\omega)$ уравнение $(\A,\omega_0)$ согласованно


(теорема \ref{th4.1}), поэтому уравнение (\ref{sys3}) имеет решение


$Z(\cdot)$, удовлетворяющее условиям $Z(0)=0$, $Z(\vt)=G$, при


$u_G(t,\omega_0)=\sum\limits_{i,p=1}^n h_{ip}u_{ip}(t,\omega_0)$, где


$u_{ip}(t,\omega_0)$ определены равенствами (\ref{4.2}), (\ref{4.3}), а


$h_{ip}$ находятся из уравнений (\ref{4.6}). Поскольку имеет место


оценка $|h|\le|\Gamma^{-1}(\vt,\omega_0)|\,|q(\vt,\omega_0)|$, а


$|\Gamma^{-1}(\vt,\omega_0)|\le\ve^{-1}$ для всех


$\omega_0\in\overline\gamma(\omega)$, то $|h|\le\ve^{-1}|q(\vt,\omega_0)|$.


Далее, найдется константа $\eta$ такая, что $|u_{ip}(t,\omega_0)|\le\eta$


при всех $(t,\omega_0)\in[0,\vt]\times\overline\gamma(\omega)$ и всех $i,p=1,\dotsc,n$


(это следует из ограниченности $\A$ на $\overline\gamma(\omega)$ и


инвариантности $\overline\gamma(\omega)$ относительно $f^t$), поэтому


найдется такое $\beta>0$, что $|u_G(t,\omega_0)|\le\beta|G(\omega_0)|$,


$(t,\omega_0)\in[0,\vt]\times\overline\gamma(\omega)$.


\end{pf}





\begin{thm}


\label{th4.3}


Пусть $\overline\gamma(\omega)$ минимально. Тогда уравнение $(\A,\omega)$


равномерно согласованно в том и только том случае, если оно согласованно.


\end{thm}





Доказательство основано на рассуждениях работы \cite{PT94}.





\begin{exam}


\label{ex1}


Система уравнений


$$


\Dot x_1=u(a_{11}(t)x_1+a_{12}(t)x_2),\quad


\Dot x_2=u(a_{21}(t)x_1+a_{22}(t)x_2),\quad |u|\le 1,


$$


с почти периодическими в смысле Бора коэффициентами равномерно


согласованна (и, следовательно, равномерно локально достижима), если


функции $a_{ij}(t)$ линейно независимы на $\R$ (тогда в силу почти


периодичности они будут линейно независимы на $[\tau,\tau+\vt]$ при


некотором $\vt>0$ и всех $\tau$). Для доказательства этого факта


достаточно по матрице $A(t)=\{a_{ij}(t)\}$ построить замыкание (в


топологии равномерной сходимости на прямой) множества сдвигов $\cal H(A)$


матрицы $A(t)$ и на $\cal H(A)$ задать динамическую систему сдвигов


(\cite{NS}, гл.\,VI, \S\,5). Так как $\cal H(A)$ минимально, то,


воспользовавшись утверждением в) теоремы \ref{th4.1} и теоремой


\ref{th4.3}, получаем требуемое.


\end{exam}





Введем в рассмотрение отображение $\Vc:M_n\to\R^{n^2}$, которое


``разворачивает'' матрицу $H=\{h_{ij}\}_{i,j=1}^n$ по строкам в


вектор-столбец


$\Vc H\doteq\col(h_{11},\dotsc,h_{1n},\dotsc,h_{n1},\dotsc,h_{nn})$.


Нетрудно проверить, что для любых $C,L,A,M\in M_n$ из $C=LA$ следует


$\Vc C=(L\otimes I_n)\Vc A$ и из $C=AM$ следует


$\Vc C=(I_n\otimes M^*)\Vc A$ ($\otimes$ --- кронекерово (прямое)


произведение матриц (\cite{Lan}, с.\,235)). Построим по уравнению


$(\A,\omega)$ так называемое ``большое уравнение'' \cite{PT95}


\begin{equation}


\label{5.1}


\Dot z=F(f^t\omega)z+G(f^t\omega)v,\quad z\in\R^{n^2},\quad v\in\R^r,


\end{equation}


где $F(\omega)=A_0(\omega)\otimes I-I\otimes A_0^*(\omega)\in M_{n^2}$,


$G(\omega)=(\Vc A_1(\omega),\dotsc,\Vc A_r(\omega))\in M_{n^2\times r}$. Через


$Z_0(t,s,\omega)$ обозначим функцию Коши уравнения $\Dot z=F(f^t\omega)z$.


Уравнение (\ref{5.1}) будем отождествлять с тройкой $(F,G,\omega)$.





Напомним, что уравнение $(F,G,\omega)$ называется вполне управляемым


(\cite{K}, гл.\,6, \S\,20), если при некотором $\vt>0$ матрица управляемости


\begin{equation}


\label{W}


W(\vt,\omega)=\int_0^{\vt}Z_0(0,t,\omega)G(f^t\omega)G^*(f^t\omega)Z_0^*(0,t,\omega)\,dt


\end{equation}


определенно положительна, и равномерно вполне управляемым \cite{T83},


если найдется такое $\vt>0$, что $W(\vt,\omega_0)$ определенно положительна


равномерно относительно $\omega_0\in\overline\gamma(\omega)$.





\begin{thm}


\label{th5.1}


Уравнение $(\A,\omega)$ согласованно $($равномерно согласованно$)$


в том и только том случае, если уравнение $(F,G,\omega)$ вполне управляемо


$($равномерно вполне управляемо$)$.


\end{thm}





\begin{pf} Пусть $\Gamma(\vt,\omega)$ --- матрица согласования


уравнения $(\A,\omega)$. Поскольку


$\wh A_l(t,\omega)=X_0(0,t,\omega)A_l(f^t\omega)X_0(t,0,\omega)$, то


$\Vc\wh A_l(t,\omega)=(X_0(0,t,\omega)\otimes X_0^*(t,0,\omega))\Vc A_l(f^t\omega)$,


поэтому $\Gamma(\vt,\omega)=\int\limits_0^{\vt}\sum\limits_{l=1}^r


\Vc\wh A_l(t,\omega)\cdot(\Vc\wh A_l(t,\omega))^*\,dt$. Далее нетрудно


проверить, что $Z_0(0,t,\omega)=X_0(0,t,\omega)\otimes X_0^*(t,0,\omega)$,


поэтому


$Z_0(0,t,\omega)G(f^t\omega)=(\Vc\wh A_1(t,\omega),\dotsc,\Vc\wh A_r(t,\omega))$.


Следовательно, $W(\vt,\omega)=\int\limits_0^{\vt}\sum\limits_{l=1}^r \Vc\wh


A_l(t,\omega)\cdot(\Vc\wh A_l(t,\omega))^*\,dt$. Таким образом,


$\Gamma(\vt,\omega)=W(\vt,\omega)$.


\end{pf}





\section{Согласованность стационарного уравнения}





Рассмотрим уравнение (\ref{1}) с постоянной матрицей $\A(\omega)\equiv\A$


и соответствующее ему уравнение $(F,G,\omega)\equiv(F,G)$. Для каждого


$\nu=1,\dotsc,r$ введем матрицы $L_0^{\nu}=A_{\nu}$,


$L_k^{\nu}=A_0L_{k-1}^{\nu}-L_{k-1}^{\nu}A_0$, $k=1,\dotsc,n^2-1$.





\begin{thm}


\label{th5.2}


Уравнение $\A$ согласованно в том и только том случае, если


среди матриц $L_k^{\nu}$, $k=0,\dotsc,n^2-1$, $\nu=1,\dotsc,r$,


существуют $n^2$ линейно независимых.


\end{thm}





\begin{pf} Уравнение $\A$ согласованно тогда и только


тогда, когда уравнение $(F,G)$ вполне управляемо, а это эквивалентно


равенству $\rank[G,FG,\dotsc,F^{n^2-1}G]=n^2$. Применив к векторам-столбцам


матрицы $G$ отображение $\Vc^{-1}$, получим $r$ матриц


$L_0^1,\dotsc,L_0^r$. Проделаем ту же процедуру с матрицами $F^kG$,


$k=1,\dotsc,n^2-1$. Тогда на каждом шаге получим $r$ матриц


$L_k^1,\dotsc,L_k^r$. Поэтому равенство $\rank[G,FG,\dotsc,F^{n^2-1}G]=n^2$


имеет место только в том случае, если линейная оболочка матриц


$L_k^{\nu}$, $k=0,\dotsc,n^2-1$, $\nu=1,\dotsc,r$, совпадает с $M_n$.


\end{pf}





\begin{cor}


\label{cor1}


Пусть $A_0$ коммутирует с матрицами $A_1,\dotsc,A_r$. Тогда


уравнение $\A$ согласованно в том и только том случае, если линейная


оболочка матриц $A_1,\dotsc,A_r$ совпадает с $M_n$.


\end{cor}





\begin{cor}


\label{cor2}


Пусть матрицы $A_0,\dotsc,A_r$ квазитреугольные,


т.\,е.~$A_{\nu}=\{B_{ij}^{\nu}\}_{i,j=1}^s$, $B_{ii}^{\nu}\in M_{n_i}$,


$\sum\limits_{i=1}^sn_i=n$, $B_{ij}^{\nu}=0$ при  $i>j$. Тогда уравнение


$\A=(A_0,\dotsc,A_r)$ не является согласованным.


\end{cor}





\begin{pf} Для всех $\nu=1,\dotsc,r$, $k=0,\dotsc,n^2-1$


матрицы $L_k^{\nu}$ имеют такой же вид, как и матрицы $A_{\nu}$,


поэтому их линейная оболочка не совпадает с $M_n$.


\end{pf}





\begin{exam}


\label{ex2}


Покажем, что уравнение $\A=(A_0,A_1)$ не является согласованным.


Достаточно построить матрицу $W_0\in M_n$, удовлетворяющую равенствам


$(\Vc W_0)^*F^kb=0$ для всех $k=0,\dotsc,n^2-1$, где


$F=A_0\otimes I-I\otimes A_0^*$, $b=\Vc A_1$. Нетрудно проверить, что


$(\Vc W_0)^*F^kb=b^*(\Vc W_k)$, $k=0,\dotsc,n^2-1$, где


$W_k=A_0^*W_{k-1}-W_{k-1}A_0^*$. Если $\Spp A_1=0$, то при $W_0=I$


выполнено $W_k=0$, $k=1,\dotsc,n^2-1$ и $b^*(\Vc W_0)=0$.


Если $\Spp A_1\ne 0$ и $A_0\ne\mu I$, то положим


$W_0=A_0^*-\dfrac{\Spp(A_1A_0)}{\Spp A_1}I$. В этом случае также выполнены


равенства $b^*(\Vc W_k)=0$ для всех $k=0,\dotsc,n^2-1$. Если


$A_0=\mu I$, то $F=0$ и согласованности нет.


\end{exam}





\begin{cor}


\label{cor3}


Если $\Spp A_{\nu}=0$ для всех ${\nu}=1,\dotsc,r$, то уравнение $\A$ не


является согласованным.


\end{cor}





{\bf Доказательство.} Положим $W_0=I$, тогда


$(\Vc W_0)^*[G,FG,\dotsc,F^{n^2-1}G]=0$.





\begin{exam}


\label{ex3}


Пусть степень минимального многочлена матрицы $A_0$ равна $m$ ($m\ge 2$).


Тогда уравнение $\A\equiv(A_0,\dotsc,A_r)$, где $r\le m-1$, не является


согласованным. Для доказательства этого достаточно построить матрицу


$W_0\in M_n$, удовлетворяющую равенствам $b_{\nu}^*(\Vc W_k)=0$ при


всех $\nu=1,\dotsc,r$, $k=0,\dotsc,n^2-1$, где $b_{\nu}=\Vc A_{\nu}$,


$W_k=A_0^*W_{k-1}-W_{k-1}A_0^*$. Будем искать $W_0$ в виде


$W_0=\sum\limits_{s=0}^{m-1}\lambda_s(A_0^*)^s$ (матрицы $(A_0^*)^s$,


$s=0,\dotsc,m-1$, линейно независимы). Тогда $W_k=0$ для всех


$k=1,\dotsc,n^2-1$. Далее, равенства $b_{\nu}^*(\Vc W_0)=0$ эквивалентны


равенствам


$\sum\limits_{s=0}^{m-1}\lambda_s\Spp(A_{\nu}A_0^s)=0$, $\nu=1,\dotsc,r$,


образующим систему из $r$ уравнений с $m$


неизвестными $\lambda_0,\dotsc,\lambda_{m-1}$. Так как $r<m$, то существует


нетривиальное решение системы, что и доказывает отсутствие


согласованности.


\end{exam}





\begin{cor}


\label{cor4}


Если все собственные значения матрицы $A_0$ различны, то


уравнение $\A\equiv(A_0,\dotsc,A_r)$, где $r\le n-1$, не является


согласованным.


\end{cor}





Таким образом, если минимальный многочлен матрицы $A_0$ совпадает с


характеристическим, то $r$ управляющих параметров, где $r<n$,


недостаточно для согласованности. Покажем, что $n$ параметров


управления достаточно, т.~е. найдутся такие матрицы $A_1,\dotsc,A_n$,


что уравнение $\A=(A_0,A_1,\dotsc,A_n)$ согласованно. Пусть


$\lambda^n+\alpha_1\lambda^{n-1}+\dotsb+\alpha_n$ --- характеристический многочлен


матрицы $A_0$. Тогда матрица $A_0$ в некотором


базисе имеет вид


$$


A_0=


\begin{pmatrix}


0 & 1 & 0 &\dots & 0\\


0 & 0 & 1 &\dots & 0\\


\vdots & \vdots & \vdots & \ddots & \vdots\\


0 & 0 & 0 & \dots & 1\\


-\alpha_n & -\alpha_{n-1} & -\alpha_{n-2} & \dots & -\alpha_1


\end{pmatrix}


.


$$


Построим матрицы $A_1=e_n e_1^*,\dotsc,A_n=e_ne_n^*$ и по уравнению


$\A$ --- уравнение $(F,G)$. Тогда


$$


F=


\begin{pmatrix}


-A_0^* & I & 0 & \dots & 0\\


0 & -A_0^* & I & \dots & 0\\


\vdots & \vdots & \vdots & \ddots & \vdots\\


0 & 0 & 0 & \dots & I\\


-\alpha_nI & -\alpha_{n-1}I & -\alpha_{n-2}I & \dots & -\alpha_1I-A_0^*


\end{pmatrix}


,\qquad


G=


\begin{pmatrix}


0\\ 0\\ \vdots\\ 0\\ I


\end{pmatrix}


.


$$


Очевидно, что $\rank[G,FG,\dotsc,F^{n-1}G]=n^2$, поэтому уравнение $\A$


согласованно.








\begin{thebibliography}{25}





\bibitem{L}


Лейхтвейс К. {\em Выпуклые множества}. -- М.: Наука, 1985.~-- 335~с.





\bibitem{M69S}


Миллионщиков В.М. {\em Доказательство достижимости центральных


показателей линейных систем} // Сиб. матем. журн. -- 1969. --


Т.\,10. -- \No\,1. -- С.\,99--104.





\bibitem{M69D}


Миллионщиков В.М. {\em Грубые свойства линейных систем


дифференциальных уравнений} // Дифференц. уравнения. -- 1969. --


Т.\,5. -- \No\,10. -- С.\,1775--1784.





\bibitem{Izob}


Изобов Н.А. {\em Линейные системы дифференциальных уравнений} //


Итоги науки и техн.\linebreak ВИНИТИ. Матем. анализ. -- 1974. -- Т.\,12. --


С.\,71--146.





\bibitem{PT94}


Попова С.Н., Тонков Е.Л. {\em Управление показателями Ляпунова


согласованных систем}. I // Дифференц.  уравнения. -- 1994. --


Т.\,30. -- \No\,10. -- С.\,1687--1696.





\bibitem{PT95}


Попова С.Н., Тонков Е.Л. {\em К вопросу о равномерной


согласованности линейных систем} // Дифференц. уравнения. --


1995. -- Т.\,31. -- \No\,4. -- С.\,723--724.





\bibitem{T95}


Тонков Е.Л. {\em Задачи управления показателями Ляпунова} //


Дифференц. уравнения. -- 1995. -- Т.\,31. -- \No\,10. -- С.\,1682--1686.





\bibitem{PT97}


Попова С.Н., Тонков Е.Л. {\em Согласованные системы и управление


показателями Ляпунова} // Дифференц. уравнения. -- 1997. -- Т.\,33. --


\No\,2. -- С.\,226--235.





\bibitem{Er}


Еругин Н.П. {\em Линейные системы обыкновенных дифференциальных


уравнений}. -- Минск: Изд-во АН БССР, 1963. -- 272 с.





\bibitem{T74}


Тонков Е.Л. {\em Управляемость нелинейной системы по линейному


приближению} // ПММ. -- 1974. -- Т.\,38. -- \No\,4. --


С.\,599--606.





\bibitem{T82}


Тонков Е.Л. {\em Равномерная локальная управляемость и стабилизация


нелинейной рекуррентной системы} // Дифференц.  уравнения. -- 1982. --


Т.\,18. -- \No\,5. -- С.\,908--910.





\bibitem{NS}


Немыцкий В.В., Степанов В.В. {\em Качественная теория


дифференциальных уравнений}. -- М.: ГИТТЛ, 1949. -- 550 с.





\bibitem{Lan}


Ланкастер П. {\em Теория матриц}. -- М.: Наука, 1978. -- 280 с.





\bibitem{K}


Красовский Н.Н. {\em Теория управления движением}. -- М.: Наука,


1968. -- 475 с.





\bibitem{T83}


Тонков Е.Л. {\em О равномерной локальной управляемости линейного


уравнения} // Матем. физика. Респ. межвед. сб. -- Киев, 1983. --


Вып.\,33. -- С.\,44--53.





\end{thebibliography}





\vskip 2cm





\post{\it Удмуртский государственный}{\it Поступила}


\post{\it университет}{03.03.1998}





\label{end}


\end{document}





